\chapter{Introduction}

Financial institutions process large volumes of customer transactions every day. While most of these transactions reflect legitimate economic activity, a small proportion may be associated with financial crime. One important example is the use of money mule accounts. Money mules are individuals or entities whose accounts are used to receive and transfer illicit funds on behalf of criminal networks. These accounts may be opened deliberately for criminal purposes, but they may also belong to individuals who have been recruited, deceived, or coerced into moving money. Detecting such accounts is therefore an important component of anti-money laundering and fraud prevention systems.

Money mule detection is challenging for several reasons. First, mule accounts are rare relative to the overall customer population, creating a highly imbalanced detection problem. Second, mule behaviour may not be captured by a single transaction or a single account-level statistic. Instead, suspicious behaviour often arises from combinations of transaction patterns, such as incoming funds being rapidly forwarded onwards, one incoming transfer being followed by multiple outgoing payments, or unusually high outgoing intensity shortly after receipt of funds. Third, many legitimate customers may display superficially similar behaviour. For example, active retail customers may make frequent outgoing payments, receive salary or family transfers, or spend incoming funds over time. This makes it difficult to distinguish mule-like behaviour from normal customer activity using simple rules alone.

This project investigates whether account-level behavioural features can be used to identify customer accounts with mule-like transaction patterns. The analysis was conducted using bank transaction data aggregated at the party level, so that each observation represents the behaviour of one customer rather than one individual transaction. This choice reflects the practical objective of identifying suspicious customer accounts for further review. A set of behavioural features was engineered to describe transaction activity, outgoing intensity, pass-through behaviour, one-in-many-out patterns, draining behaviour, and activity timing. These features were designed to capture both general account activity and more specific mule-related behaviours.

The modelling strategy combines unsupervised and supervised approaches. DBSCAN was used as the main unsupervised clustering method to segment customer behaviour and identify observations that do not belong to dense regions of normal retail activity. An autoencoder was also used as an unsupervised reconstruction-based anomaly detection model, where accounts with large reconstruction errors were interpreted as more unusual relative to the training distribution. These unsupervised methods are useful because confirmed mule labels are limited and the majority of accounts are unlabeled.

In addition, a small set of confirmed mule accounts was used to construct supervised baseline models. Logistic regression was used as an interpretable linear baseline, while XGBoost was used as a more flexible non-linear tree-based model capable of capturing interactions between behavioural features. Since confirmed mule accounts formed only a small minority of the modelling dataset, class weighting was applied to reduce the tendency of supervised models to favour the majority unlabeled class. The supervised model outputs were interpreted as relative mule-risk scores rather than calibrated probabilities, because the confirmed mule accounts were appended to a sampled retail comparison group and the unlabeled class may contain hidden mule accounts.

The main objective of this thesis is to develop and evaluate a behavioural detection framework for identifying mule-like accounts in bank transaction data. The project combines customer-level feature engineering with both unsupervised anomaly detection and supervised classification models, so that suspicious accounts can be segmented, scored, and ranked according to mule-risk. In particular, the thesis compares DBSCAN and autoencoder-based anomaly detection with logistic regression and XGBoost supervised baselines, evaluates the extent to which engineered behavioural features improve detection, and analyses which transaction patterns drive the resulting risk scores under a partially labelled setting.

The remainder of this thesis is organised as follows. Chapter 2 provides background on money mule behaviour, anomaly detection, and supervised learning methods relevant to financial crime detection. Chapter 3 describes the data and preprocessing steps. Chapter 4 presents the feature engineering process, including activity-based features and mule-behaviour features. Chapter 5 describes the modelling methodology, including DBSCAN, the autoencoder, logistic regression, and XGBoost. Chapter 6 presents and interprets the experimental results. Finally, Chapter 7 discusses the limitations of the study and outlines possible directions for future work.




